\input{../tex-inputs/latex-leseansicht-vorspann}

\section[Gustav Schwarzkopf an Arthur Schnitzler, {[}10.? 3. 1903{]}]{L04331 Gustav Schwarzkopf an Arthur Schnitzler, {[}10.? 3. 1903{]}}
\nopagebreak\mylabel{L04331v}
\rehead{ }\normalsize\beginnumbering\briefempfaengerindex{Schnitzler, Arthur@\textsc{Schnitzler, Arthur}!zzzSchwarzkopf, Gustav@\emph{von Gustav Schwarzkopf}!1903-03-101@{{[}10.? 3. 1903{]}}|(be}
\toendnotes[C]{\smallbreak\pagebreak[2]}
\correspDesc{Versand  durch Gustav Schwarzkopf am [10.? 3. 1903] in Wien
\newline{}Erhalt  durch Arthur Schnitzler am [10.? 3. 1903] in Wien}\toendnotes[C]{\smallbreak}
\Standort{CUL, Schnitzler, A 117,4.}
\physDesc{Brief, Durchschlag, 5 Blätter, 5 Seiten, 3849 Zeichen
\newline{}Schreibmaschine
\newline{}HandschriftX2 einer Schreibkraft: Bleistift (\noindent{}marginale Korrekturen)
\newline{}Schnitzler: mit Bleistift Vermerk: »März 1903, nach der (Vor-) Lecture
                                    des urſprüngl Doppelſtücks\pwindex{Schnitzler, Arthur 15. 5. 1862 Wien – 21. 10. 1931 ebd.@\textsc{Schnitzler, Arthur} (15. 5. 1862 Wien – 21. 10. 1931 ebd.), \emph{Schriftsteller, Mediziner}!einsame Weg. Schauspiel in fünf Akten@\strich\emph{Der einsame Weg. Schauspiel in fünf Akten}|pwv}\pwindex{Schnitzler, Arthur 15. 5. 1862 Wien – 21. 10. 1931 ebd.@\textsc{Schnitzler, Arthur} (15. 5. 1862 Wien – 21. 10. 1931 ebd.), \emph{Schriftsteller, Mediziner}!Professor Bernhardi. Komödie in fünf Akten@\strich\emph{Professor Bernhardi. Komödie in fünf Akten}|pwv} von \uline{\textsc{G.
                                          Schwarzkopf}}« }
\buchAbdrucke{\weitereDrucke{Arthur Schnitzler: \emph{Der einsame Weg. Historisch-kritische Ausgabe}. Herausgegeben von Anna Lindner und Isabella Schwentner unter Mitarbeit von Teresa Klestorfer, Marina Rauchenbacher und Miriam Sibitz. Berlin, Boston: \emph{deGruyter} 2024, S. 1236–1238.} }\toendnotes[C]{\smallbreak}
\pstart
           \noindent{}{\pb}\label{K_L04318-1v}\edtext{Johanna ist nicht die Tochter}{\lemma{\textnormal{\emph{Johanna … Tochter}}}\Cendnote{\textnormal{Am 19. 2. 1903 hatte Schnitzler{ }Schwarzkopf\pwindex{Schwarzkopf, Gustav 7.\,11.\,1853 Wien – 13.\,11.\,1939 ebd.@\textsc{Schwarzkopf, Gustav} (7.\,11.\,1853 Wien – 13.\,11.\,1939 ebd.), \emph{Schriftsteller}|pwk} und Olga Gussmann\pwindex{Schnitzler, Olga 17.\,1.\,1882 Wien – 13.\,1.\,1970 Lugano@\textsc{Schnitzler, Olga} (17.\,1.\,1882 Wien – 13.\,1.\,1970 Lugano), \emph{Schauspielerin, Sängerin}|pwk} den Stand der Arbeit am »Doppelstück\pwindex{Schnitzler, Arthur 15. 5. 1862 Wien – 21. 10. 1931 ebd.@\textsc{Schnitzler, Arthur} (15. 5. 1862 Wien – 21. 10. 1931 ebd.), \emph{Schriftsteller, Mediziner}!einsame Weg. Schauspiel in fünf Akten@\strich\emph{Der einsame Weg. Schauspiel in fünf Akten}|pwkv}\pwindex{Schnitzler, Arthur 15. 5. 1862 Wien – 21. 10. 1931 ebd.@\textsc{Schnitzler, Arthur} (15. 5. 1862 Wien – 21. 10. 1931 ebd.), \emph{Schriftsteller, Mediziner}!Professor Bernhardi. Komödie in fünf Akten@\strich\emph{Professor Bernhardi. Komödie in fünf Akten}|pwkv}« vorgelesen. Im
                     \emph{Tagebuch}\pwindex{Schnitzler, Arthur 15. 5. 1862 Wien – 21. 10. 1931 ebd.@\textsc{Schnitzler, Arthur} (15. 5. 1862 Wien – 21. 10. 1931 ebd.), \emph{Schriftsteller, Mediziner}!Tagebuch@\strich\emph{Tagebuch}|pwk} hielt er dazu das Urteil fest:
                     »Ganz wie ich vermuthete: 2 Stücke\pwindex{Schnitzler, Arthur 15. 5. 1862 Wien – 21. 10. 1931 ebd.@\textsc{Schnitzler, Arthur} (15. 5. 1862 Wien – 21. 10. 1931 ebd.), \emph{Schriftsteller, Mediziner}!einsame Weg. Schauspiel in fünf Akten@\strich\emph{Der einsame Weg. Schauspiel in fünf Akten}|pwv}\pwindex{Schnitzler, Arthur 15. 5. 1862 Wien – 21. 10. 1931 ebd.@\textsc{Schnitzler, Arthur} (15. 5. 1862 Wien – 21. 10. 1931 ebd.), \emph{Schriftsteller, Mediziner}!Professor Bernhardi. Komödie in fünf Akten@\strich\emph{Professor Bernhardi. Komödie in fünf Akten}|pwv}; in jetziger Fassung
                     unmöglich.« Schwarzkopf\pwindex{Schwarzkopf, Gustav 7.\,11.\,1853 Wien – 13.\,11.\,1939 ebd.@\textsc{Schwarzkopf, Gustav} (7.\,11.\,1853 Wien – 13.\,11.\,1939 ebd.), \emph{Schriftsteller}|pwk}
                  verfasste in den Folgetagen ein Szenario, wie er glaubte, dass sich die Probleme
                  lösen ließen. Wann und in welcher Form Schnitzler den Text von Schwarzkopf\pwindex{Schwarzkopf, Gustav 7.\,11.\,1853 Wien – 13.\,11.\,1939 ebd.@\textsc{Schwarzkopf, Gustav} (7.\,11.\,1853 Wien – 13.\,11.\,1939 ebd.), \emph{Schriftsteller}|pwk} erhielt, ist nicht mit Sicherheit zu ermitteln. 
                  Das vorliegende Typoskript
                  dürfte aber von Schnitzler beauftragt
                  worden sein und nicht das Original darstellen. Durch die fehlenden Bezugnahmen auf den
                  Text im \emph{Tagebuch}\pwindex{Schnitzler, Arthur 15. 5. 1862 Wien – 21. 10. 1931 ebd.@\textsc{Schnitzler, Arthur} (15. 5. 1862 Wien – 21. 10. 1931 ebd.), \emph{Schriftsteller, Mediziner}!Tagebuch@\strich\emph{Tagebuch}|pwk} und in der Korrespondenz 
                  ist es wahrscheinlich, dass Schwarzkopf\pwindex{Schwarzkopf, Gustav 7.\,11.\,1853 Wien – 13.\,11.\,1939 ebd.@\textsc{Schwarzkopf, Gustav} (7.\,11.\,1853 Wien – 13.\,11.\,1939 ebd.), \emph{Schriftsteller}|pwk} sein Manuskript 
                  persönlich übergab, wofür sich das erste Wiedersehen der beiden am A. S.: \emph{Tagebuch}, 10. 3. 1903 nach einer Reise Schnitzlers
                  anbietet.}}}\label{K_L04318-1}, sondern nur eine entfernte
               Verwandte, die im Hause erzogen wurde. Zwischen ihr und Hans eine Art Liebe, mit Ehe
               in Aussicht, die aus Kameradschaft, gemeinsamen Interessen, Ansichten entstanden
               ist.\pend
           
\pstart
           Eine jüngere Freundin der Frau Pflugfelder |:in der Nachbarschaft wohnend:|. Ihr Mann
               ist einige Wochen nach der Hochzeit unheilbar erkrankt. Sie pflegt ihn seit fünfzehn
               Jahren.\pend
           
\pstart
           Sie hat die Einleitungsszene mit Johanna, die erklärt, dass sie nicht leiden sehen
               kann, sich wundert, dass Frau Anna ihren Mann nicht verlassen hat. Sie hat ja auch
               nur ein Leben, das ebensoviel wert sei wie das seine.\pend
           
\pstart
           Frau Anna: sie liebe ihn wol nicht mehr, sie empöre sich wol auch zuweilen gegen das
               Schicksal, aber sie habe nun einmal dieses Los gezogen – Pflicht – sie hindert ihn,
               sich zu töten. Sala hat ihr einmal vor langen Jahren von Liebe gesprochen. Sie hat
               ihn zurückgewiesen.\pend
           
\pstart
           {\pb}Dann die Szenenführung wie bisher. Prof.
               Bernhardi, das Gespräch des Hans mit ihm. Nur müssten die Schrecken, die der
               erkrankten Frau harren, in der Unterredung mit Bernhardi stärker betont werden,
               ebenso wie das ganz ungewöhnlich innige Verhältnis zwischen Mutter und Sohn, das die
               scherzhaft geäusserte und ernst gemeinte Eifersucht des Vaters hervorgerufen hat.\pend
           {\vspace{1\baselineskip}}
\pstart
           II. Akt\pend
           
\pstart
           Unmittelbar nach dem Begräbnis. Man erfährt, wie die Frau einige Wochen gelitten.\pend
           
\pstart
           Der Mann ganz gebrochen.\pend
           
\pstart
           Johanna zu ihrer Ueberraschung tiefer erschüttert als sie glaubte.\pend
           
\pstart
           Hans in tiefstem Schmerz, aber gefasst, nur irritiert von dem wortreichen Schmerz der
               Andern. Verwundert, dass Niemand ausspricht, dass der schnelle Ausgang doch
               eigentlich noch ein Glück sei.\pend
           
\pstart
           {\pb}Bernhardi ahnt, da er Hansens Ansichten kennt,
               und giebt es ihm zu verstehen.\pend
           
\pstart
           Dann irgend ein Kondolenzbesuch von Frauen.\pend
           
\pstart
           Zum ersten Mal fällt das Wort »Erlösung«. Hans klammert sich daran. Mit dem ganz
               vernichteten Vater allein geblieben, greift er das Wort wieder auf, verallgemeinert,
               ob man nicht dürfte {\dots} man müsste{\dots} Dann erzählt er. Er hat seine Mutter überrascht, wie sie sich zu dem Schränkchen
               tastete, zu dem der Vater und er die Schlüssel haben. Er hat sie gefragt und da hat
               sie ihn auf den Knieen gebeten, ein Ende zu machen. Er brauche es nicht selbst zu
               tun: ein Versehen, er brauche nur den Schlüssel stecken zu lassen. Das hat er
               getan.\pend
           
\pstart
           Dann wie bisher, nur das Entsetzen, der Abscheu des Vaters stärker.\pend
           {\vspace{1\baselineskip}}
\pstart
           {\pb}III. Akt.\pend
           
\pstart
           Einige Tage später.\pend
           
\pstart
           Hans, Johanna. Auch sie wendet sich von ihm ab, sie empfindet Grauen in seiner Nähe.
               Er erinnert sie, wie stark sie sich gegeben, dass sie seine Ansichten geteilt.\pend
           
\pstart
           Sie sagt, sie sei eben nur ein ganz gewöhnliches Weib, ihre Geistesstärke war nur
               Anempfindung, nur der Wunsch, ihm zu gefallen. Jetzt kann sie’s garnicht denken, dass
               sie ihm angehören könnte.\pend
           
\pstart
           Es kommt die Nachricht, dass Frau Anna’s Mann gestorben sei. Sie kommt vielleicht
               selbst und erzählt, dass er in schmerzfreien Stunden doch noch gern gelebt habe. Ob
               er die nicht beneide?\pend
           
\pstart
           Dann ein paar Füllszenen.\pend
           
\pstart
           Etwa Sala |:auch Arzt:|, der seine Reise aufgeben will, um Anna zur Seite zu
               stehen.\pend
           
\pstart
           Schlussszene zwischen Vater und Sohn. Hans bereut nicht, fühlt sich nicht schuldig,
               spricht {\pb}für seine Ueberzeugung. Er durfte so
               handeln, weil er sie am besten, am wahrsten geliebt hat, – darum musste er ihr diesen
               letzten Liebesdienst erweisen. Wie, wenn Einer an seiner Ehre unheilbar erkrankt ist,
               drückt Ihr ihm den Revolver in die Hand. Besser der Tod als die Schande. Das
               bewundert Ihr, das gebt Ihr zu, das ist ein Liebesbeweis. Warum denn nicht auch
               besser der Tod als entsetzliches Leiden?\pend
           
\pstart
           Er geht fort, für Sala, nach Asien. Nicht um den Tod zu suchen, nicht um zu büssen,
               nur um den Personen, die seinen Anblick nicht ertragen können, aus dem Wege zu gehen.
               Und wenn er doch zurückkommt – und er wird sich tapfer wehren –, wird ihm sein Vater
               vielleicht doch wieder die Hand drücken können.\pend
           
\pstart
           \raggedleft{}Gustav Schwarzkopf.\pend
           \selectlanguage{ngerman}\endnumbering\briefempfaengerindex{Schnitzler, Arthur@\textsc{Schnitzler, Arthur}!zzzSchwarzkopf, Gustav@\emph{von Gustav Schwarzkopf}!1903-03-101@{{[}10.? 3. 1903{]}}|)be}\mylabel{L04331h}
\begin{anhang}
\end{anhang}\newcommand{\dateiname}{L04331}\newcommand{\titel}{Gustav Schwarzkopf an Arthur Schnitzler, [10.? 3. 1903]}\newcommand{\editorInnen}{Herausgegeben von Jahnke, SelmaMüller, Martin Anton}\input{../tex-inputs/latex-leseansicht-abspann}
